\section{A valued lift and descent of marked covers}\label{sec:valuations}

Fix the group $J$, the integer $h\geq 1$, and $n=p^s$ from
\Cref{thm:tuple-divisibility}, and retain the characteristic-$p$ data
$A_0,P,F',F_0,E,t_1,\ldots,t_n$ of \Cref{prop:twist-index}.
Thus $P$ has exact order $n$, $F'=\kappa(A_0^n)$, and the generator
of $\operatorname{Gal}(F'/F_0)$ satisfies
$\sigma(t_i)=t_{i+1}-P$, with cyclic indices.  All the valued fields
constructed in this section are auxiliary; they are independent of the
block coefficient system $(K,\mathcal O,k)$.

The output sought here is a cover over a field of small $p$-part degree
whenever the finite marked-cover set has a small Sylow orbit. We first
lift the curve and its labels, then choose a valued base on which full
field degrees survive as residue degrees. This prepares the genus
comparison in \Cref{sec:specialization} without yet constructing a
model of the cover.

\subsection{Lifting the twist and its marked points}

\begin{proposition}\label{prop:mixed-lift}
There are complete discretely valued fields $K_b\subset K_b'$ of
characteristic zero, with valuation rings $R_b\subset R_b'$, and a smooth
projective genus-one curve $\mathcal X/R_b$ with the following properties.
\begin{enumerate}[label=\textup{(\roman*)}]
\item The extension $K_b'/K_b$ is cyclic Galois of degree $n$;
the extension $R_b'/R_b$ is finite \emph{\'{e}tale} Galois, its residue
extension is $F'/F_0$, and the two fields have the same value group.
Moreover, $K_b$ contains $\mathbb Q_p$ and all $m$th roots of unity,
where $m=|J|$.
\item The special fiber of $\mathcal X$ is $E/F_0$.
There is a relatively ample effective divisor $D$ on $\mathcal X$ of
fiber degree $n$.
\item Over $R_b'$ there are pairwise disjoint sections
$\xi_1,\ldots,\xi_n$ reducing respectively to $t_1,\ldots,t_n$.
They avoid $D$, and their union descends to a finite \emph{\'{e}tale}
effective divisor on $\mathcal X$.
Under the cyclic descent action the sections are permuted by
$\xi_i\mapsto\xi_{i+1}$.
\end{enumerate}
\end{proposition}

\begin{proof}
We first construct the constant complete DVR.  Starting with
$\mathbb Z_p$, successively lift monic separable irreducible polynomials
over the residue fields so that the resulting finite residue fields
exhaust $\kappa=\overline{\mathbb F}_p$.  To justify this construction,
if $R$ is the current complete DVR and $f\in R[T]$ is such a monic
lift, then $R[T]/(f)$ is finite free, its reduction is a field, and its
maximal ideal is generated by the original uniformizer.  The polynomial
$f$ is irreducible over the fraction field: a monic factorization there
would have coefficients integral over $R$, hence in $R$, and would give
a factorization of its irreducible reduction.  Consequently the quotient
is a local domain with principal maximal ideal, hence a DVR; its residue
extension is the prescribed one.  Its derivative is a unit, so the
extension is unramified.  The union of this tower is a valuation ring
with value group $\mathbb Z$ and residue $\kappa$.  Complete it.  Completion
preserves its uniformizer and residue field, giving a complete DVR of
mixed characteristic with residue $\kappa$. Denote it for the moment by
$R_c^{(0)}$, with fraction field $K_c^{(0)}$.

Lift the coefficients of the smooth Weierstrass equation for $A_0$ to
this DVR.  Its discriminant remains a unit, so it defines a smooth
elliptic scheme $\mathcal A$ by \Cref{foundation:elliptic}.
We next lift $P$ without assuming that the torsion scheme is \emph{\'{e}tale}.
The kernel $\mathcal A[n]$ is finite locally free of degree $n^2$.
Its generic fiber is dense: in its finite flat coordinate algebra the
uniformizer is a nonzerodivisor, so no irreducible component is supported
in the special fiber. Since this dense generic fiber is a finite set,
one of its points $Q$ has closure containing $P$. Its field of definition
$K_Q$ is finite over $K_c^{(0)}$; this is a characteristic-zero field,
not a valuation residue field. The reduced closure of $Q$ is finite and
integral over $R_c^{(0)}$. Pass to the integral closure of $R_c^{(0)}$
in $K_Q$.
By \Cref{foundation:valuation} it is a complete DVR, and lying over
shows that its closed point maps to $P$.  We obtain an $n$-torsion section
$\widetilde P$ specializing to $P$.  Its order is exactly $n$, since its
order divides $n$ and reduction has order $n$.  Enlarge the constant
field once more to contain all $m$th roots of unity.  These are finite
extensions of complete discretely valued fields, and their residue
fields are still $\kappa$, which is algebraically closed.  Denote the
resulting constant valuation ring by $R_c$ and choose its uniformizer
$\varpi$.

Consider the smooth product $\mathcal A^n/R_c$.  Its generic fiber is
integral, and flatness makes the total space integral.  At the generic
point of the integral special fiber, its local ring is a noetherian local
domain whose maximal ideal is generated by $\varpi$ and whose residue
field is $F'$.  It is therefore a DVR.  Complete its fraction field for
this valuation and call the result $K_b'$, with valuation ring $R_b'$.
The product projections give points $\xi_i\in\mathcal A(K_b')$.
Their sections over $R_b'$ reduce to the generic projections $t_i$.

The automorphism of the product specified by
\[
             \sigma(\xi_i)=\xi_{i+1}-\widetilde P
\]
fixes the constants, has order $n$, and preserves the special fiber.
It preserves the valuation and extends to $K_b'$.  Its residue action
is the specified action on $F'$, of order $n$.  Put
$K_b=(K_b')^{\langle\sigma\rangle}$.  The fixed field is closed in the
complete valued field $K_b'$, so it is complete.  The fixed uniformizer
$\varpi$ shows that its value group is the same discrete group.
Finite fixed-field theory gives
$[K_b':K_b]=n$ with Galois group $\langle\sigma\rangle$.

Let $k_b$ be the residue field of $K_b$.  Then $k_b\subseteq F_0$.
Lifting a linearly independent system in a residue extension to units
gives a linearly independent system over the valued base field:
normalize a putative relation by a coefficient of minimum value and
reduce.  Applying this observation here gives
\[
 n=[F':F_0]\leq [F':k_b]\leq[K_b':K_b]=n.
\]
Thus $k_b=F_0$.

To verify the stronger assertion about valuation rings, choose a
primitive element $\bar a$ of the separable extension $F'/F_0$ and a lift
$a\in R_b'$.  The polynomial
\[
              f(T)=\prod_{j=0}^{n-1}(T-\sigma^j a)\in R_b[T]
\]
reduces to the separable irreducible polynomial of $\bar a$.
Consequently $R_b[a]$ is finite free of rank $n$, local with maximal
ideal $(\varpi)$, and a domain with fraction field $K_b'$.
It is a DVR, and the derivative $f'(a)$ is a unit.  Uniqueness of the
extension of the complete-DVR valuation identifies it with $R_b'$.
This proves that $R_b'/R_b$ is finite \emph{\'{e}tale}.
The Galois torsor map
\[
 R_b'\otimes_{R_b}R_b'\longrightarrow\prod_{j=0}^{n-1}R_b',
 \qquad a\otimes b\longmapsto (a\sigma^j(b))_j
\]
is an isomorphism: after reduction it is the usual isomorphism for
$F'/F_0$, and both sides are free $R_b'$-modules of rank $n$.

On $\mathcal A_{R_b'}$ define the semilinear descent action
$\delta(x)=\sigma(x)+\widetilde P$. Here $\sigma$ acts on the
coefficients, whereas $\delta$ also translates the elliptic-curve point.
Its $n$th power is the identity.
The divisor
\[
                    D'=\sum_{j=0}^{n-1}[j\widetilde P]
\]
is invariant.  Its sections are disjoint because their reductions are
distinct.  The line bundle $\mathcal O(3D')$ is a tensor product of
translates of the relatively ample Weierstrass line bundle
$\mathcal O(3[0])$.  It is relatively ample, and therefore so is
$\mathcal O(D')$.  The invariant divisor supplies its canonical
compatible linearization.  Finite \emph{\'{e}tale} Galois descent
(\Cref{foundation:descent}) gives $\mathcal X$ and $D$, with a splitting
identification $\mathcal X_{R_b'}\simeq\mathcal A_{R_b'}$.
Their reduction is precisely the twist $E$ and the corresponding
origin-orbit divisor.  Smoothness and relative ampleness descend.

Under this identification, continue to write $\xi_i$ for the corresponding
sections of $\mathcal X_{R_b'}$. They have distinct reductions, all different
from the constant points $jP$.  They are therefore pairwise disjoint and
avoid $D'$. Under the descent action they satisfy
$\delta(\xi_i)=\xi_{i+1}$. Their union is an invariant
disjoint union of sections, hence a finite \emph{\'{e}tale} divisor of
degree $n$, and descends with these properties, although the individual
sections need not be defined over $R_b$.
\end{proof}

\subsection{A spherical base and its finite extensions}

For a valued field we write the valuation additively, with
$v(0)=+\infty$.  A \emph{closed ball} of radius $\delta$ means
$B(a,\delta)=\{x:v(x-a)\geq\delta\}$, where $\delta$ belongs to the
value group.  A valued field is \emph{spherically complete} if every
nested family of nonempty closed balls has nonempty intersection.
An extension is \emph{immediate} if its value group and residue field
are unchanged.

The purpose of the next construction is twofold. Finite field degrees
will equal full residue degrees, so a small-orbit degree bound survives
reduction. Finite-dimensional section spaces will also have free
integral lattices with full-dimensional reduction, allowing comparison
of the generic and residue genera. The construction uses the classical
maximal-immediate-extension and pseudo-limit methods of Krull and
Kaplansky \cite{Krull1932,Kaplansky1942}; see
\cite[Section~5]{Poonen1993} for an accompanying account. We give the
details for our specified residue field, without assuming it perfect
or algebraically closed.

\begin{proposition}\label{prop:spherical-base}
There is an extension $K_s/K_b$ with value group $\mathbb Q$ and residue
field $F_0$ which is spherically complete.  Every finite extension
$L/K_s$, equipped with an extension of its valuation, has value group
$\mathbb Q$, is spherically complete, and satisfies
\begin{equation}\label{eq:full-residue-degree}
                   [\operatorname{res}(L):F_0]=[L:K_s].
\end{equation}
The valuation on $K_s$, and on each such $L$, has a unique extension to
every algebraic extension.  Moreover
$\overline{K_s}\simeq\mathbb C$ as abstract fields.
\end{proposition}

We prove the construction and the linear-algebra assertions first, and
establish uniqueness after proving simultaneous residue approximation.

\begin{proof}[Construction of $K_s$]
Normalize the valuation of \Cref{prop:mixed-lift} by $v(\varpi)=1$.
In an algebraic closure with an extended valuation, choose compatible
factorial roots: take $\theta_1=\varpi$ and
$\theta_{j+1}^{j+1}=\theta_j$ for $j\geq1$. Thus
$\theta_j^{j!}=\varpi$ and the fields $K_b(\theta_j)$ are nested.
Put $K_\infty=\bigcup_{j\geq1}K_b(\theta_j)$.
For $r=j!$ and $\theta=\theta_j$, one has $\theta^r=\varpi$.
The values of $1,\theta,\ldots,\theta^{r-1}$ lie in distinct cosets of
$\mathbb Z$.  These elements are linearly independent over $K_b$, and
the equation for $\theta$ shows that they are a basis of $K_b(\theta)$.
For $x=\sum_{j=0}^{r-1}a_j\theta^j\ne0$, the distinct value cosets give
\[
                  v(x)=\min_j\bigl(v(a_j)+j/r\bigr).
\]
The value group is $r^{-1}\mathbb Z$.  If $v(x)=0$, the minimum can only
occur at $j=0$, and all other terms have positive value; hence the residue
of $x$ is the residue of $a_0$.  This proves that $K_\infty$ has value
group $\mathbb Q$ and residue field $F_0$.

Here is the set-theoretic detail needed to choose a maximal immediate
extension.  Any valued field with countable value group and countable
residue field has cardinality at most $\mathfrak c=2^{\aleph_0}$.
Indeed, for every $q$ in the value group and every closed $q$-ball $B$,
the open $q$-balls contained in $B$ are in bijection with the residue
field, after choosing a center and an element of value $q$.
Label these open subballs injectively by natural numbers.
For a point $x$, record, for each $q$, the label of its open $q$-ball
inside its closed $q$-ball.  If $x\ne y$, at $q=v(x-y)$ they belong to
the same closed ball but different open subballs, so these records
differ.  This gives an injection into a countable product of countable
sets.  In our situation $F_0$ is countable, being a subfield of the
finitely generated field $F'/\kappa$.

Choose a fixed set of cardinality greater than $\mathfrak c$, containing
the underlying set of $K_\infty$.  Order the immediate valued extensions
of $K_\infty$ carried by subsets of this set by literal extension of all
their structures.  The union of a chain is an immediate valued field:
each of its elements and every finite field computation occur in one
member of the chain.  The cardinality bound applies to the union as
well.  Zorn's lemma gives a maximal member $K_s$.  Any further immediate
extension has cardinality at most $\mathfrak c$ and could be relabeled,
fixing $K_s$, into the unused part of the fixed set.  Thus $K_s$ has no
proper immediate extension in the unrestricted sense.

We now prove spherical completeness directly.  Suppose a nest of
nonempty closed balls has empty intersection.  Distinct balls in a nest
with the same radius coincide, so the countability of the value group
gives a countable cofinal subnest.  Successively go deeper in that
subnest, also choosing each next ball to omit the preceding center.
We obtain cofinal balls $B(a_i,\delta_i)$ such that
$a_i\notin B(a_{i+1},\delta_{i+1})$ and
\begin{equation}\label{eq:pseudoconvergent-centers}
 v(a_j-a_i)=\gamma_i\quad(j>i),\qquad
 \delta_i\leq\gamma_i<\delta_{i+1},\qquad
 \gamma_1<\gamma_2<\cdots .
\end{equation}
The balls $B(a_i,\gamma_i)$ are nested, lie inside $B(a_i,\delta_i)$,
and still have empty intersection.

Extend the valuation to an algebraic closure of $K_s$.
Call $b$ a \emph{pseudo-limit} of this sequence if
$v(b-a_i)=\gamma_i$ for all sufficiently large $i$.
No element of $K_s$ is a pseudo-limit, since a pseudo-limit belongs to
every ball $B(a_i,\gamma_i)$.
For any point $b$ which is not a pseudo-limit, the values $v(b-a_i)$
are eventually constant and smaller than $\gamma_i$.
To see this, if $v(b-a_i)<\gamma_i$ at one stage, the strong triangle
law and \eqref{eq:pseudoconvergent-centers} give the same value at every
later stage.  If $v(b-a_i)>\gamma_i$, the next value is $\gamma_i$, which
is smaller than $\gamma_{i+1}$, reducing to the preceding case.
Otherwise equality holds eventually, as required for a pseudo-limit.

Suppose first that an algebraic pseudo-limit $\alpha$ exists, and
choose one of smallest degree $d$ over $K_s$.  Every root $\beta$ of a
nonzero polynomial $f\in K_s[T]$ with $\deg f<d$ is not a pseudo-limit.
For large $i$ we therefore have
\[
 v(a_i-\beta)=v(\alpha-\beta)<v(\alpha-a_i),\qquad
 \frac{\alpha-\beta}{a_i-\beta}\equiv1
       \pmod{\mathfrak m_{\overline{K_s}}}.
\]
Factoring $f$ over the algebraic closure yields
\begin{equation}\label{eq:pseudolimit-leading-term}
             f(\alpha)/f(a_i)\equiv1
                    \pmod{\mathfrak m_{\overline{K_s}}}
             \qquad(i\gg0).
\end{equation}
Every nonzero element of $K_s(\alpha)$ is $f(\alpha)$ for a polynomial
of degree less than $d$.  Equation \eqref{eq:pseudolimit-leading-term}
therefore gives, for each such element $x$, a scalar $c=f(a_i)\in K_s$
with $v(x-c)>v(x)$. It follows that $v(c)=v(x)$, so no new values
occur. If $v(x)=0$, the same inequality says that $x$ and $c$ have
the same residue, so no new residue-field elements occur either.
Thus $K_s(\alpha)/K_s$ is a proper immediate extension, a contradiction.

Suppose instead that no algebraic pseudo-limit exists.  Factoring any
nonzero polynomial over the fixed algebraic closure shows that its
values at $a_i$ are eventually constant and finite; the same holds for
nonzero rational functions.  Define, on $K_s(T)$,
\[
                 w(f)=\text{the eventual value of }v(f(a_i)),
                 \qquad w(0)=+\infty.
\]
Evaluation proves multiplicativity and the strong triangle inequality,
so $w$ is a valuation extending that of $K_s$.
For fixed $f\ne0$ and $j>i\gg0$, every linear factor in its numerator
and denominator satisfies
$(a_j-\beta)/(a_i-\beta)\equiv1\pmod{\mathfrak m_{\overline{K_s}}}$,
where $\beta$ is its root. Indeed $v(a_i-\beta)<\gamma_i=v(a_j-a_i)$.
Consequently
\[
                    v(f(a_j)-f(a_i))>v(f(a_i))=w(f).
\]
For a fixed sufficiently large $i$, the left side has an eventual
value as $j$ increases, since it evaluates the rational function
$f(T)-f(a_i)$.  Thus
$w(f-f(a_i))>w(f)$ unless the difference is zero, when the assertion
is immediate. Thus every nonzero element $f$ of $K_s(T)$ admits a
scalar $f(a_i)$ with this strict approximation inequality. As in the
algebraic case, this shows that neither the value group nor the residue
field changes. This is again a proper immediate extension,
contradicting maximality.  Both cases being impossible, $K_s$ is
spherically complete.
\end{proof}

\begin{lemma}\label{lem:orthogonal-bases}
Let $K_*$ be a spherically complete field with value group $\mathbb Q$,
valuation ring $R_*$, maximal ideal $\mathfrak m_*$, and residue field
$k_*$.  Let $V$ be a finite-dimensional $K_*$-vector space with a
valuation norm $\lambda:V\to\mathbb Q\cup\{+\infty\}$: it is finite
away from zero, has $\lambda(0)=+\infty$, satisfies the strong triangle inequality, and obeys
$\lambda(ax)=v(a)+\lambda(x)$ for $a\ne0$.
There is a basis $e_1,\ldots,e_d$ of norm zero such that
\begin{equation}\label{eq:orthogonal-basis}
             \lambda\left(\sum_i a_i e_i\right)=\min_i v(a_i).
\end{equation}
In particular its integral and positive balls are
\begin{equation}\label{eq:orthogonal-lattice}
 V_{\geq0}=\bigoplus_i R_*e_i,\qquad
 V_{>0}=\bigoplus_i\mathfrak m_*e_i
             =\mathfrak m_*V_{\geq0},
 \qquad \dim_{k_*}(V_{\geq0}/V_{>0})=d.
\end{equation}
The normed space $V$ is spherically complete.
\end{lemma}

\begin{proof}
A subspace with an orthogonal basis is spherically complete: its closed
balls are products of closed balls in the finitely many coordinates,
with the radii shifted by the norms of the basis vectors.  Each
coordinate nest has a common point.

Suppose an orthogonal basis has already been constructed for a subspace
$U$, and choose $x\notin U$.  For every value $\gamma$ achieved by
$\lambda(x-u)$, $u\in U$, consider
\[
                  B_\gamma=\{u\in U:\lambda(x-u)\geq\gamma\}.
\]
This is a nonempty closed ball in $U$: if $u_\gamma$ belongs to it,
then it equals $\{u:\lambda(u-u_\gamma)\geq\gamma\}$.
These balls form a nest, so they have a common point $u_0$.
The finite value $\ell=\lambda(x-u_0)$ is at least every achieved value
and itself is achieved.  Thus $u_0$ is a best approximation to $x$.
Put $y=x-u_0$.  For $u\in U$, maximality of $\ell$ gives
$\lambda(y+u)\leq\ell$.  If $\lambda(u)<\ell$, the strong triangle
law gives $\lambda(y+u)=\lambda(u)$; if $\lambda(u)\geq\ell$, it gives
$\lambda(y+u)\geq\ell$, hence equality.  Therefore
\[
                       \lambda(y+u)=\min\{\ell,\lambda(u)\}.
\]
Homogeneity shows that adjoining $y$ preserves orthogonality.
Induction gives an orthogonal basis of $V$.  Since every basis norm is
in the base value group, scale its vectors to norm zero.
Formulae \eqref{eq:orthogonal-basis} and \eqref{eq:orthogonal-lattice}
follow, and the same coordinate argument proves spherical completeness
of $V$.
\end{proof}

\begin{proof}[Finite-extension assertions of \Cref{prop:spherical-base}]
For any algebraic extension of a field with value group $\mathbb Q$,
every nonzero algebraic element has value in $\mathbb Q$.
Indeed, in a polynomial relation $\sum_i a_i x^i=0$, at least two
nonzero terms must have the same minimum value, so
\[
                       (i-j)v(x)=v(a_j)-v(a_i)\in\mathbb Q
                       \quad(i\ne j).
\]
Divisibility of $\mathbb Q$ and torsion-freeness of ordered value groups
imply $v(x)\in\mathbb Q$.
Apply \Cref{lem:orthogonal-bases} to a finite extension $L/K_s$ with any
extended valuation, viewed as a normed vector space over $K_s$.
Its integral and positive balls are exactly its valuation ring and
maximal ideal.  Formula \eqref{eq:orthogonal-lattice} proves
\eqref{eq:full-residue-degree}, including the full inseparable degree
of the residue extension.  Orthogonal coordinates also show that $L$
is spherically complete.  The same argument applies to every finite
extension of $L$.
\end{proof}

The preceding argument gives equality of field and residue degrees for
finite extensions of the spherical base. We also need a statement for
valued fields that are not assumed spherically complete, where several
prolongations can occur. The next lemma supplies the residue-degree bound
used for curve function fields in \Cref{sec:specialization} and completes
the uniqueness assertion for the spherical base.

\begin{lemma}[Simultaneous residue approximation]\label{lem:simultaneous-residues}
Let $T$ be a nontrivially valued field with value group $\mathbb Q$ and
infinite residue field $k_T$.  Let $U/T$ be finite and let
$w_1,\ldots,w_r$ be distinct prolongations of its valuation, with residue
fields $k_1,\ldots,k_r$.  Then
\[
             \bigcap_{i=1}^r\mathcal O_{w_i}\longrightarrow
                         \prod_{i=1}^r k_i
\]
is surjective, and
\begin{equation}\label{eq:residue-degree-bound}
                         \sum_{i=1}^r[k_i:k_T]\leq[U:T].
\end{equation}
\end{lemma}

\begin{proof}
The preceding polynomial-relation argument puts all the values in
$\mathbb Q$, with their normalization fixed by the valuation on $T$.
For $i\ne j$, distinctness gives an element whose two values differ.
Inverting it if necessary and multiplying by an element of $T$ of a
suitable rational value produces $z\in U$ with
$w_i(z)>0>w_j(z)$.
Choose a unit $c\in T$ such that, at every $w_\ell$ for which
$w_\ell(z)=0$, the residue of $1+cz$ is nonzero.
Each such condition excludes at most one value for the residue of $c$
in $k_T^\times$, so the infinitude of $k_T$ permits this choice.
Then
\[
                         q_{ij}=\frac{1}{1+cz}
\]
is integral at every $w_\ell$, has residue $1$ at $i$, and has residue
$0$ at $j$.  Indeed, a positive value of $z$ makes the denominator a
unit with residue $1$; a negative value of $z$ gives the denominator
that negative value; and value zero is covered by the choice of $c$.

Put $q_i=\prod_{j\ne i}q_{ij}$, with $q_i=1$ if $r=1$.
It has residue $1$ at $i$ and positive value at all other places.
For a prescribed residue $b_i\in k_i$, choose any lift $\widetilde b_i$
integral at $i$. If the lift is nonzero, choose $N_i$ so that
\[
 N_iw_\ell(q_i)+w_\ell(\widetilde b_i)>0
 \qquad(\ell\ne i).
\]
There are only finitely many conditions, and every $w_\ell(q_i)$ in
them is positive. At $i$, the factor $q_i^{N_i}$ has residue one.
Thus $q_i^{N_i}\widetilde b_i$ is integral everywhere, has residue
$b_i$ at $i$, and has residue zero at every other place. A zero lift
contributes the zero term. Summing these terms realizes any tuple $(b_i)_i$.

For completeness, each $[k_i:k_T]$ is finite and at most $[U:T]$ by the
single-place lifting argument already used in the proof of
\Cref{prop:mixed-lift}.  Choose a $k_T$-basis of each $k_i$, and use the
surjectivity just proved to lift its elements with zero residues in the
other coordinates.  These lifts are linearly independent over $T$.
For if they satisfied a nonzero relation, divide it by a coefficient
of minimum value.  Every coefficient then is integral and at least one
is a unit.  Reducing in a coordinate which contains such a unit
contradicts the chosen residue basis there.  Their total number is
therefore at most $[U:T]$, proving \eqref{eq:residue-degree-bound}.
\end{proof}

\begin{proof}[Completion of the proof of \Cref{prop:spherical-base}]
The residue of $K_s$, and that of every finite extension $L/K_s$, is
infinite.  If two distinct valuations extended the valuation of $L$ to
a finite extension $U/L$, the finite-extension argument above would
give residue degree $[U:L]$ for each.  This contradicts
\eqref{eq:residue-degree-bound}.  Existence of prolongations is part of
\Cref{foundation:valuation}; uniqueness on finite extensions implies
uniqueness on arbitrary algebraic extensions by restriction to finite
subextensions.

Finally $|K_s|\leq\mathfrak c$ by the cardinality argument, while
$\mathbb Q_p\subset K_s$ gives $|K_s|\geq\mathfrak c$.
Its algebraic closure has the same cardinality.
An uncountable algebraically closed field of characteristic zero has
cardinality equal to its transcendence degree over $\mathbb Q$: the
field generated by a transcendence basis, and then its algebraic
closure, both have cardinality the maximum of that basis cardinality
and $\aleph_0$.  Both $\overline{K_s}$ and $\mathbb C$ therefore have
transcendence degree $\mathfrak c$ over $\mathbb Q$.
The algebraically closed field classification in
\Cref{foundation:fields} gives the asserted abstract isomorphism.
No compatibility with either field's topology is asserted or used.
\end{proof}

\subsection{Marked covers and a small orbit}

Choose an embedding $K_b'\hookrightarrow\overline{K_s}$ over $K_b$ and
an abstract isomorphism $\overline{K_s}\simeq\mathbb C$.
Write $X_s=\mathcal X_{K_s}$. The embedding carries the sections
$\xi_i$ on the split model $\mathcal X_{R_b'}$ to the specified
geometric marked points on $X_s$.
A \emph{marked $J$-cover} of $\mathcal X_{\overline{K_s}}$ will mean a connected
smooth projective curve $Y$ with a finite map to $\mathcal X_{\overline{K_s}}$
which is Galois of group $J$, together with a specified left action of
$J$ by deck transformations.  Isomorphisms are over the fixed base
curve and commute with every element of this specified action.

\begin{lemma}\label{lem:marked-cover-classification}
After choosing oriented handle and puncture generators on the complex
punctured curve, isomorphism classes of marked $J$-covers branched only at
$\xi_1,\ldots,\xi_n$ correspond bijectively to
\[
 \left\{(x,y,u_1,\ldots,u_n)\in J^{n+2}:
 \begin{array}{l}
 [x,y]u_1\cdots u_n=1,\\
 \langle x,y,u_1,\ldots,u_n\rangle=J
 \end{array}\right\}/J,
\]
where $J$ acts by simultaneous inner conjugation.
The inertia order at $\xi_i$ is $|u_i|$.
For a prescribed multiset $\mathcal C$ of $n$ $J$-conjugacy classes of
$p$-singular elements, let $\mathcal T_{\mathcal C}$ be the finite set
of these marked-cover classes having that multiset of local classes.
Then $\mathcal T_{\mathcal C}$ is preserved by
$\operatorname{Gal}(\overline{K_s}/K_s)$.
\end{lemma}

Thus $|\mathcal T_{\mathcal C}|=N_{J,s}(\mathcal C)$ in the notation
of \Cref{thm:tuple-divisibility}. The labels $\xi_1,\ldots,\xi_n$
remain ordered: prescribing their class multiset does not divide out
by permutations of the labels.

\begin{proof}
Under the chosen complex identification the base is a smooth
Weierstrass cubic. In characteristic zero, completing the square gives
an equation $y^2=f_3(x)$ with three distinct roots. Its double cover of
the projective line branches simply at those roots and at infinity.
Cut the sphere along two disjoint arcs joining pairs of these four
points and glue the two sheets crosswise. The resulting compact surface
is a torus. This identifies the topological type of our particular
base without a general analytic--algebraic genus comparison.

Remove the $n$ marked points. Its complex fundamental group has
generators and relation
$[x,y]u_1\cdots u_n=1$ as in \Cref{foundation:complex}.
Choose a point in a fiber of a marked cover and identify the fiber with
the regular left $J$-set.  Path lifting commutes with the left deck
action, so monodromy acts by right translations.  With path multiplication
chosen in the order of successive traversal, these translations specify
a homomorphism from the fundamental group to $J$.  The cover is
connected exactly when the monodromy action is transitive, which for
this regular action means that the homomorphism is surjective.
Changing the chosen fiber point conjugates this homomorphism by one
element of $J$.  Conversely, a surjective homomorphism defines such a
fiber action and therefore a connected unramified cover.
An equivariant fiber bijection is a right translation, and compatibility
with monodromy says exactly that the two homomorphisms are simultaneously
conjugate.  Thus the quotient is by inner conjugation; keeping the
specified $J$-action does not identify covers through outer automorphisms
of $J$.

Riemann existence in \Cref{foundation:complex} algebraizes the
unramified cover of this punctured algebraic curve, with its morphisms
and specified deck action. Normalize the original projective base in
the resulting function field. Finiteness of normalization gives a
projective algebraic cover; normality makes its source smooth in
characteristic zero. The deck action and isomorphisms extend uniquely
across normalization. In a local power-map chart $z\mapsto z^e$, the
ramification index is $e$, the order of the associated monodromy
element. These extensions give precisely the branched covers and
isomorphisms in the statement.

We describe the local conjugacy class algebraically in order to check
Galois invariance.  For each divisor $e$ of $m$, use the root
$\zeta_e\in K_s$ which maps to $\exp(2\pi i/e)$ under the chosen
complex identification.  All these roots belong to $K_s$ by
\Cref{prop:mixed-lift}.
Represent the based puncture loop as an approach path, a positive
circle around the puncture, and the reverse approach path.
In the chart $z\mapsto z^e$, lifting the positive circle sends $z$ to
$\zeta_e z$. Let $Q$ be the point above the puncture selected by the
lifted approach path, and let $g\in J$ be the inertia element acting
this way in its local chart. Apply $g$ to the lifted approach path.
The reverse of the resulting path is the lift of the return path
starting at the endpoint of the circle lift. Thus the complete based
loop sends the chosen fiber point to its image under $g$.
Under the preceding fiber identification, its monodromy element is
therefore $g$. The inertia group is cyclic, and $g$ is its unique
generator acting on the tangent line at $Q$ by $\zeta_e$.
Choosing another point above the puncture conjugates the element in $J$.
This recovers its conjugacy class from the algebraic tangent action.
(Reversing all local-loop conventions replaces all these generators by
their inverses consistently.)

A Galois automorphism over $K_s$ sends the tangent action at $Q$ to
the tangent action at its conjugate point, and fixes $\zeta_e$.
Consequently it preserves the local $J$-conjugacy class at the transported
branch point.  It permutes the branch points, since their union is
defined over $K_s$, and hence preserves their prescribed class multiset
and the condition that all inertia orders are divisible by $p$.
This proves invariance of $\mathcal T_{\mathcal C}$.
Its finiteness follows from the finite tuple description.
Although Galois can permute the marked points, an isomorphism of the
marked covers is over the fixed geometric base curve. The cover classes
are therefore still parametrized by simultaneous inner conjugacy,
without quotienting by those permutations.
\end{proof}

For a cover with its group action specified, the obstruction to descent
is center-valued; compare the marked-cover analysis of
D\`ebes and Douai \cite[Section~3.2]{DebesDouai1997}. In the following
proposition we realize that obstruction as an explicit finite group
extension and split it by a Sylow argument. Thus a field of moduli is
not silently assumed to be a field of definition.

\begin{proposition}[Descent from a small Sylow orbit]\label{prop:small-orbit-descent}
Assume $O_p(J)=1$ and fix a multiset $\mathcal C$ as in
\Cref{lem:marked-cover-classification}.  There is a finite Galois
extension $M/K_s$ through which the action on $\mathcal T_{\mathcal C}$
factors.  Let $P_\Gamma$ be a Sylow $p$-subgroup of
$\Gamma=\operatorname{Gal}(M/K_s)$.
If a $P_\Gamma$-orbit has size less than $p^h$, one of its marked covers
descends to a finite extension $L/K_s$ satisfying
\begin{equation}\label{eq:small-orbit-degree}
                [L:K_s]_p<p^h,
                \qquad [F:F_0]_p=[L:K_s]_p,
                \qquad F=\operatorname{res}(L).
\end{equation}
The descended map $Y\to X=\mathcal X_L$ has geometrically connected
smooth projective source, degree $m$, the specified $J$-action, and all
branch indices at the $\xi_i$ divisible by $p$.
\end{proposition}

\begin{proof}
Choose a representative for each of the finitely many marked-cover
classes.  Their projective equations, maps, and specified deck actions
are finite data, and hence are defined over one finite extension of
$K_s$ by \Cref{foundation:descent}.  Take a finite Galois extension
$N/K_s$ containing these fields of definition, and enlarge it if necessary
to contain $K_sK_b'$.  For every $\tau\in\operatorname{Gal}(N/K_s)$
and every chosen representative, choose a marked isomorphism between
its $\tau$-conjugate and the representative of the resulting class.
Such an isomorphism exists over $\overline{K_s}$ by
\Cref{lem:marked-cover-classification}; there are only finitely many
choices to make.  Their finite coefficient data lie in a further
finite extension, which we enclose in a finite Galois extension $M/K_s$.
This two-stage choice ensures that all required isomorphisms are
defined over $M$.  Conjugates of the original models under
$\operatorname{Gal}(M/K_s)$ depend only on restriction to $N$ and are
scalar extensions of the already considered conjugates.  The action
on classes therefore factors through $\Gamma$.

Let $S_p\leq P_\Gamma$ stabilize a class in a small orbit and let $Y_M$
be its chosen model. Since $S_p$ fixes this isomorphism class, for each
$\tau\in S_p$ the chosen isomorphisms supply a $\tau$-semilinear
isomorphism of $Y_M$ over the corresponding action on $\mathcal X_M$,
commuting with the specified $J$-action. Here semilinear means that the
map changes scalars by $\tau$ rather than fixing $M$. These maps have
not been chosen to respect multiplication in $S_p$; that compatibility
is the remaining descent problem.

Form the group $\mathcal E$ of all such marked semilinear isomorphisms,
for all $\tau\in S_p$, with composition as its group operation.
Projection to the scalar action is a surjection $\mathcal E\to S_p$.
An automorphism of $Y_M$ over $\mathcal X_M$ is a deck transformation: after
extension to $\overline{K_s}$ the Galois cover has exactly its $m$
specified deck transformations, all already defined over $M$.
Those commuting with the specified $J$-action are exactly $Z(J)$.
Every marked semilinear map also commutes with these specified deck
transformations. Thus this kernel is central, and $\mathcal E$ is
finite and fits into an exact sequence
\begin{equation}\label{eq:semilinear-cover-extension}
                  1\longrightarrow Z(J)\longrightarrow\mathcal E
                     \longrightarrow S_p\longrightarrow1.
\end{equation}
The Sylow $p$-subgroup of the abelian group $Z(J)$ is characteristic
there, hence normal in $J$.  Since $O_p(J)=1$, it is trivial; thus
$p\nmid|Z(J)|$.
A Sylow $p$-subgroup of $\mathcal E$ has order $|S_p|$, intersects
$Z(J)$ trivially, and therefore maps isomorphically onto $S_p$.
Its inverse supplies a section of \eqref{eq:semilinear-cover-extension},
assigning maps $\phi_\tau$ with
\[
 \phi_\tau\phi_\rho=\phi_{\tau\rho},\qquad \phi_1=\mathrm{id}.
\]
These are exactly the compatibility identities for a descent datum on
$Y_M$ and its cover map. Each $\phi_\tau$ commutes with every specified
element of $J$, so the datum descends the marked action as well.

Take the pullback of an ample line bundle on $X_s$, for example
$\mathcal O(D)$ on its generic fiber.  It is ample on $Y_M$, and the
semilinear cover maps give its compatible linearization.
Galois descent in \Cref{foundation:descent} therefore gives the marked
cover over $L=M^{S_p}$.  Its claimed geometric properties can be checked
after extension to $\overline{K_s}$, where they are those of the chosen
cover.  In particular its function-field extension is Galois of degree
$m$ with the specified group $J$.
Finally
\[
 [L:K_s]_p=[\Gamma:S_p]_p
       =\frac{|P_\Gamma|}{|S_p|}
       =|P_\Gamma:S_p|<p^h.
\]
Equation \eqref{eq:full-residue-degree} gives the residue assertion in
\eqref{eq:small-orbit-degree}.
\end{proof}

\subsection{Matching branch permutations after reduction}

In \Cref{sec:specialization} we will impose vanishing over selected marked
points. The same subset of label indices must define a branch subset over
the characteristic-zero field and over its residue field. We therefore
compare their Galois permutation actions. These are the actions on the
marked curve; in split elliptic-curve coordinates they include the
translating torsion point in the descent data.

\begin{lemma}\label{lem:residue-permutations}
Let $L/K_s$ be finite, with residue $F$, and set
$L'=LK_b'$.  In the residue field of $L'$, form $S=FF'$ using the
residues of its two subfields.  Then $L'/L$ and $S/F$ are finite Galois
extensions, the residue field of $L'$ is exactly $S$, and reduction
induces an isomorphism
\[
             \operatorname{Gal}(L'/L)\simeq\operatorname{Gal}(S/F).
\]
Under this isomorphism the permutation actions on $\xi_1,\ldots,\xi_n$
and on $t_1,\ldots,t_n$ agree. In particular, every subset of indices
invariant under these identified permutation groups defines a branch
subset over $L$ and over $F$, respectively.
\end{lemma}

\begin{proof}
The uniqueness proved in \Cref{prop:spherical-base} gives one valuation
on $L'$.  Its restriction to $K_b'$ is the valuation of
\Cref{prop:mixed-lift}, by uniqueness over the complete discretely
valued field $K_b$.  Put $T=\operatorname{res}(L')$.
It contains both $F$ and $F'$, so it contains their compositum $S$.
Base change of the finite Galois extension $K_b'/K_b$ gives the Galois
extension $L'/L$, and restriction embeds
$H=\operatorname{Gal}(L'/L)$ in $\operatorname{Gal}(K_b'/K_b)$.
Every element of $H$ preserves the unique valuation and hence acts on
$T$, fixing $F$ and preserving $F'$.  Its induced action on $S$ is
faithful: if it is trivial on $S$, it is trivial on $F'$; the faithful
residue action of $\operatorname{Gal}(K_b'/K_b)$ then makes its
restriction to $K_b'$ trivial; and $L$ and $K_b'$ generate $L'$.
Thus reduction injects $H$ into $\operatorname{Aut}_F(S)$.
Also $S/F$ is Galois, being a compositum with the finite Galois
extension $F'/F_0$.  The residue-degree bound now gives the complete
degree comparison
\begin{equation}\label{eq:residue-permutation-sandwich}
 |H|\leq[S:F]\leq[T:F]\leq[L':L]=|H|.
\end{equation}
All these inequalities are equalities.  In particular $T=S$ and the
injection of groups is an isomorphism.

The sections $\xi_i$ extend over the valuation ring of $L'$ by base
change from the split good-reduction model, and reduce to $t_i$.
Their reductions are distinct.  Applying a member of $H$ to a section
and then reducing gives the same section as first reducing and then
applying its image in $\operatorname{Gal}(S/F)$.  This proves equality
of the two permutations of the index set.  An invariant subset of a
split finite \emph{\'{e}tale} branch divisor descends along the respective
finite Galois extension.  The final assertion follows by taking the
same invariant subset of indices under these identified permutation
actions: it selects the $\xi_i$ on the generic side and the $t_i$ on
the residue side.
\end{proof}

\Cref{fig:field-tower} records these inclusions and residue fields in
the situation of \Cref{prop:small-orbit-descent}.
\begin{figure}[!ht]
\centering
\begin{tikzpicture}[
  every node/.style={font=\small},
  field/.style={inner sep=4pt},
  inclusion/.style={-{Stealth[length=4pt]},semithick}
]
\begin{scope}[xshift=-3.45cm]
  \node[font=\small\bfseries] at (0,4.65) {Characteristic zero};
  \node[field] (kb) at (0,0) {$K_b$};
  \node[field] (ks) at (-1.45,1.35) {$K_s$};
  \node[field] (kbp) at (1.45,1.35) {$K_b'$};
  \node[field] (l) at (-1.45,2.8) {$L$};
  \node[field] (lp) at (0,4) {$L'=LK_b'$};
  \draw[inclusion] (kb)--(ks);
  \draw[inclusion] (kb)--(kbp);
  \node[font=\footnotesize,align=left] at (1.85,.35)
    {unramified\\degree $n$};
  \draw[inclusion] (ks)--(l);
  \draw[inclusion] (l)--(lp);
  \draw[inclusion] (kbp)--(lp);
\end{scope}
\begin{scope}[xshift=3.45cm]
  \node[font=\small\bfseries] at (0,4.65) {Residue fields};
  \node[field] (f0) at (0,0) {$F_0$};
  \node[field] (fp) at (1.45,1.35) {$F'$};
  \node[field] (f) at (-1.45,2.8) {$F$};
  \node[field] (s) at (0,4) {$S=FF'$};
  \draw[inclusion] (f0)--(fp);
  \draw[inclusion] (f0)--(f);
  \draw[inclusion] (f)--(s);
  \draw[inclusion] (fp)--(s);
\end{scope}
\node[font=\footnotesize,align=center] at (0,-.85)
 {$\operatorname{res}(K_b)=\operatorname{res}(K_s)=F_0,
   \qquad [F:F_0]=[L:K_s],\qquad [L:K_s]_p<p^h.$};
\end{tikzpicture}
\caption{The auxiliary fields in a descent arising from a small Sylow orbit. Arrows denote field inclusions. The right-hand diagram records the residues of the left-hand fields; in particular $K_b$ and $K_s$ have the same residue. The field $K_b'$ splits the marked sections, and its residue $F'$ splits their reductions. The cyclic extension $K_b'/K_b$ reduces to $F'/F_0$, and the Galois actions of $L'/L$ and $S/F$ induce the same permutations of the marked points.}
\label{fig:field-tower}
\end{figure}

